This section has two aims. First, we identify the principal limitations of the present analysis and clarify which residual systematics arise from the current modelling choices. Secondly, we outline how the framework can be extended from the controlled setting studied here towards analysis-ready redshift calibration for future LSST data.

\subsection{Current limitations} \label{sec:6.1}

The present analysis is intentionally controlled. This is valuable for isolating failure modes and testing mitigation strategies, but several ingredients remain idealised relative to a full LSST application. Taken together, these limitations define a clear hierarchy: the dominant current systematic is mismatch in the simulated colour--redshift relation, especially in the faint and high-redshift regime where source calibration is most challenging. The simplified spectroscopic selections, missing imaging-level effects, limitations of the SOM-based discretisation and conditional density estimators, and restricted nuisance parameterisations mainly affect how that mismatch is sampled, propagated, and compressed. 

\subsubsection{Controlled approximations to spectroscopic selection} \label{sec:6.1.1}

Our treatment of spectroscopic selection is simplified. The \texttt{Degradation} datasets are generated from stylised cuts in colour, magnitude, and redshift, together with stochastic sampling to mimic incomplete spectroscopic success. While these constructions bracket a plausible range of calibration samples, they do not reproduce the detailed selection logic of any individual survey. Real spectroscopic compilations depend on footprints, targeting priorities, fibre assignment, observing conditions, redshift success, and heterogeneous quality cuts, yielding effective selection functions that are typically far more structured than those considered here \citep{2016arXiv161100036D}.

Consequently, the experiment ensemble may overestimate the variation expected from any particular LSST-era multi-survey compilation, and the resulting uncertainty should be interpreted as a controlled sensitivity study rather than a forecast for a specific calibration set. Even so, the exercise remains informative: in practice, LSST calibration will require choices about how to combine bright wide-area surveys with smaller, deeper samples. Those choices act as hyper-parameters of the calibration problem, and the Dirichlet-weighted scheme provides a pragmatic route to marginalise over them rather than fixing a single ad hoc mixture.

\subsubsection{Simulation--observation mismatch in the colour--redshift relation} \label{sec:6.1.2}

The dominant limitation of the augmentation framework remains mismatch between simulated and observed galaxy populations, particularly for faint and high-redshift galaxies. Improving overlap in photometric feature space does not guarantee that the simulated catalogues reproduce the correct \emph{conditional} colour--redshift relation at fixed colours and magnitudes. This issue is already apparent in Paper~I, where discrepancies between \texttt{OpenUniverse2024} and \texttt{CosmoDC2} grow in \(u-g\) and \(g-r\) towards faint magnitudes and high redshift, and where parts of the \texttt{Application} population remain absent from the augmented \texttt{Combination} sample.

The same limitation appears in the present calibration results. The highest-redshift source bins remain the hardest to recover, and the \texttt{Zinc} configuration demonstrates that support alone is insufficient when the underlying colour--redshift relation is biased. Figure~\ref{fig:11} illustrates this directly. In both LSST-Y1 and LSST-Y10, regions of SOM space that are well anchored by spectroscopy typically show small residual offsets, with extended areas close to zeros. The largest excursions occur in weakly constrained outskirts and substructures of the manifold, where spectroscopic support becomes sparse and calibration relies more heavily on augmentation. These residuals form coherent positive and negative patches --- not cell-scale noise --- indicating systematic shifts in the local colour--redshift relation. This behaviour is clear for \texttt{Restriction}, while moving to \texttt{Combination} reduces the mismatch in some transition regions (most clearly in LSST-Y1) but leaves substantial residual structure in augmentation-dominated regions, especially in LSST-Y10 where a larger fraction of the faint, high-redshift manifold depends on augmentation. The limiting systematic is therefore not only incomplete spectroscopic coverage, but also imperfect modelling of galaxy SEDs and their redshift evolution precisely where augmented samples carry most of the calibration weight.

\begin{figure*}
    \centering
    \begin{subfigure}[b]{0.48\linewidth}
        \centering
        \includegraphics[width=\linewidth]{Figures/Y1/SOM.pdf}
        \caption{LSST Y1}
    \end{subfigure}
    \hspace{0.0\linewidth}
    \begin{subfigure}[b]{0.48\linewidth}
        \centering
        \includegraphics[width=\linewidth]{Figures/Y10/SOM.pdf}
        \caption{LSST Y10}
    \end{subfigure}
    \caption{Residual scaled biases in the mean true redshift per cell of the augmentation SOM (Section~\ref{sec:3.1}), \(\delta_{\langle z_{\mathrm{true}}\rangle}\), for simulated catalogues relative to the target \texttt{Application} population. In each depth scenario, the upper map shows the \texttt{Restriction} dataset and the lower map the \texttt{Combination} dataset under the baseline experiment. The residuals form coherent large-scale structures across the SOM manifold, indicating systematic mismatch in the colour--redshift relation. Moving from \texttt{Restriction} to \texttt{Combination} reduces the mismatch in some regions, most clearly for Y1, but substantial residual structure remains, especially for Y10. White cells indicate insufficient support.}
    \label{fig:11}
\end{figure*}

\subsubsection{Missing imaging-level systematics} \label{sec:6.1.3}

The present analysis is performed at catalogue level. While the mocks include photometric errors and spectroscopic incompleteness, they do not explicitly forward-model imaging-level systematics such as star--galaxy separation \citep{2025arXiv250105739B}, depth fluctuations \citep{2016ApJS..226...24L, 2021A&A...648A..98J, 2024MNRAS.535.2970H, 2025A&A...694A.259Y}, deblending failures \citep{2023arXiv231002079R, 2025arXiv250316680L}, masking, or crowding. These effects can alter both observed photometry and effective selection, distorting the mapping between the underlying galaxy population and the calibration sample in ways that catalogue-level perturbations cannot capture. This omission is expected to matter most for the faint galaxies that dominate LSST source samples. A forward-modelled treatment of these systematics, based on synthetic source injection, is outlined in Section~\ref{sec:6.2.2}.

\subsubsection{Limitations of the discretisation and conditional density estimators} \label{sec:6.1.4}

The SOM-based discretisation adopted here is interpretable and well suited to support diagnostics, but it is not unique. The fixed rectangular lattice and neighbourhood function can introduce spurious colour--redshift degeneracies in high-dimensional photometric spaces, and the hard assignment of galaxies to best-matching units discards information about the continuous structure of the underlying manifold. Alternative manifold-learning methods may alleviate some of these shortcomings and are discussed in Section~\ref{sec:6.2.3}.

Similar considerations apply to the conditional density estimators. The \texttt{DIR} estimator relies on local histograms within SOM clusters and can therefore become noisy or biased when occupation numbers are low, a regime that is common in sparsely sampled high-redshift regions. The \texttt{Stack} estimator, which aggregates galaxy-level photo-\(z\) posteriors, does not satisfy the conditions under which posterior stacking is known to be consistent. \citet{2021PhRvD.103h3502M} show that the stacked estimator recovers the true redshift distribution only in two limiting cases: perfectly informative data, in which each photo-\(z\) PDF reduces to a delta function at the true redshift, or a perfectly uninformative prior that already equals the truth. Neither condition is realistic for broad-band photometric surveys, but the SOM-based augmentation framework developed here pushes towards both limits simultaneously. On the data side, photo-\(z\) posteriors are inherently probabilistic and can never reduce to delta functions; however, data augmentation improves the completeness of the feature space from which the ML model learns its conditional density estimates, reducing extrapolation and yielding more reliable --- and hence more informative --- posteriors. On the prior side, the SOM-based weighted sampling used to construct the combined calibration sample (Section~\ref{sec:3}) matches the augmentation galaxies to the photometric target distribution in feature space, regularising the effective training prior towards the population of interest; residual mismatches between simulated and observed populations nonetheless persist (Section~\ref{sec:6.1.2}). The performance comparison in Section~\ref{sec:4} is consistent with this picture: \texttt{Stack} performs adequately in well-supported regions of colour space but degrades where augmentation coverage is incomplete. Neither \texttt{Stack} nor the ad hoc \texttt{Hybrid} combination carries formal optimality guarantees; both are retained to characterise their empirical behaviour, with all estimators evaluated against the \texttt{Truth} benchmark so that residual biases are diagnosed rather than assumed negligible. More robust alternatives, including conditional normalising flows and hierarchical Bayesian inference, are outlined in Section~\ref{sec:6.2.3}.

\subsubsection{Limits of low-dimensional uncertainty parameterisations} \label{sec:6.1.5}

Finally, the uncertainty products explored here highlight a limitation of standard low-dimensional nuisance models. Mean shifts and affine width rescalings can capture part of the variation across experiments, but they do not represent the full freedom of the experiment-marginalised ensemble of tomographic redshift distributions, particularly for broad, asymmetric, or high-redshift source bins \citep{2006ApJ...636...21M}. This motivates more flexible compression schemes, including basis expansions and hierarchical latent-variable representations that preserve physically relevant shape variation more faithfully than \(\Delta_m\) and \(\Gamma_m\) alone \citep{2021A&A...650A.148S, 2022MNRAS.511.2170C, 2023MNRAS.518..709Z}.

\subsection{Roadmap to LSST deployment} \label{sec:6.2}

Extending this framework to LSST requires moving from a proof-of-concept into a multi-layer calibration programme linking targeted spectroscopy, simulation-informed augmentation, imaging-level forward modelling, and cosmological inference. Some elements are realistic ingredients for early LSST analyses, while others are longer-term developments towards a fully hierarchical end-to-end treatment. The near-term priority is to deliver robust calibration products for heterogeneous LSST-era datasets \citep{2019ApJ...873..111I}; the longer-term goal is to embed those products in a more complete probabilistic framework for redshift inference.

\subsubsection{Near-term deployment: targeted spectroscopy and multi-layer calibration} \label{sec:6.2.1}

A direct near-term application of the SOM framework is to identify where spectroscopy is most lacking in galaxy colour space. As a compact representation of the observed SED manifold, the SOM singles out sparsely sampled cells, poor-overlap regions between calibration and target samples, and loci where inferred redshift information is dominated by extrapolation, thereby providing a practical diagnostic for prioritising spectroscopic effort.

LSST calibration will benefit from deeper spectroscopy targeted at faint and high-redshift galaxies \citep{2014arXiv1410.4506S, 2014arXiv1410.4498A, 2015APh....63...81N}. An important ongoing effort of this kind is the Complete Calibration of the Color-Redshift Relation (C3R2) Survey \citep{2017ApJ...841..111M, 2019ApJ...877...81M, 2020A&A...642A.192E}, which obtains deep spectroscopy in under-sampled regions of the joint galaxy colour manifold relevant to \textit{Euclid}, \textit{Roman}, and LSST; its depth is nominally matched to \textit{Euclid} but the resulting calibration sample is of broad utility for LSST photometric-redshift calibration as well. Further gains are expected from upcoming facility-scale programmes, for example via DESI \citep{2022AJ....164..207D} deep-field programmes \citep{2025arXiv251215964B, 2026AJ....171...71R}, Subaru PFS \citep{2014PASJ...66R...1T}, 4MOST \citep{2016SPIE.9910E..1NW}, WEAVE \citep{2024MNRAS.530.2688J}, MUST \citep{2024arXiv241107970Z}, and related facilities. However, near-complete spectroscopy at LSST depth is unlikely. Targeted programmes should therefore be viewed as anchors in the most weakly constrained regions of colour space, while simulation-informed augmentation remains essential for transferring this information to the much larger photometric sample. Clustering-redshift techniques, which infer the redshift distribution of a photometric sample from its angular cross-correlation with a spectroscopic reference \citep{2008ApJ...684...88N, 2020A&A...642A.200V, 2022MNRAS.510.1223G, 2023MNRAS.524.5109R}, provide a complementary and largely independent calibration cross-check. Recent Stage-III analyses have demonstrated the power of combining colour-based SOMPZ calibration with clustering-redshift constraints into a unified ensemble of redshift distributions \citep{2024MNRAS.527.2010G, 2025arXiv251023566Y, 2025arXiv250319440W}, and extending such a joint framework to the present SOM-based augmentation pipeline is a natural next step.

A realistic LSST implementation will also require more than a single SOM defined on one catalogue, since relevant information will be distributed across wide-field LSST data, Deep Drilling Fields \citep{2024ApJS..275...21G}, overlap spectroscopy, and external surveys with different depths, bandpasses, and redshift quality. A natural extension is a multi-layer framework in which separate manifolds are constructed for distinct observational regimes and linked through transfer relations \citep{2019MNRAS.489..820B}. One possible realisation is a wide-field LSST SOM, a deeper optical SOM for the Deep Drilling Fields, and additional optical--NIR manifolds where external data are available (e.g.\ \textit{Euclid} and \textit{Roman}). The calibration problem then becomes one of linking multiple observational representations of the same underlying galaxy population, rather than forcing all information into a single feature space \citep{2022ApJS..258...15E}.

In parallel, Moskowitz et al. (in preparation) are applying data-augmentation techniques directly to LSST Data Preview~1 observations, providing an early observational validation of augmentation-based calibration on real Rubin data. That work will also present a detailed like-for-like comparison between the colour--magnitude-space augmentation strategy of \citet{2024ApJ...967L...6M} and the SOM-based augmentation developed in Paper~I and extended here.

\subsubsection{Forward modelling, augmentation, and uncertainty propagation} \label{sec:6.2.2}

A central next step is to address the imaging-level systematics identified in Section~\ref{sec:6.1.3} by moving beyond catalogue-level augmentation towards a forward-modelled treatment of observational effects. Synthetic source injection \citep{2016MNRAS.457..786S, 2022ApJS..258...15E, 2025OJAp....8E..65A} provides a well-established route to characterising these effects: artificial galaxies are embedded into real or simulated images using tools such as \texttt{GalSim} \citep{2015A&C....10..121R} and propagated through the full detection, measurement, and photometric-redshift estimation pipeline \citep{2021arXiv210912060V}. By comparing injected and recovered properties, the survey transfer function---encoding depth variations, PSF-dependent selection, blending, and masking---can be learned empirically rather than imposed through catalogue-level approximations. In this way, augmented samples capture not only sparse coverage in catalogue space but also the observational distortions that determine how that space is populated, including the sample variance \citep{2012MNRAS.423..909C} and distribution incompleteness \citep{2022arXiv220708657M} that arise when imaging-level selection is neglected.

The same forward-modelled setting enables a more realistic augmentation strategy than reliance on a single simulation. Future analyses can draw on an ensemble of augmentation models, including hydrodynamical simulations \citep{2015MNRAS.446..521S, 2019ComAC...6....2N, 2025arXiv250821126S, 2025arXiv250904067C}, semi-analytic or empirical models \citep{2012NewA...17..175B, 2015ARA&A..53...51S, 2023A&A...670A.100L}, and synthetic SED catalogues spanning varied assumptions about galaxy--halo connection, star-formation history, dust, emission lines, and black-hole accretion and feedback \citep{2013ARA&A..51..393C, 2024ApJS..274...12A, 2025ApJ...993..240T, 2025A&A...703A.255T, 2025JCAP...06..007F}. Treating these as distinct augmentation hypotheses, Dirichlet-weighted sampling across model families can propagate not only calibration-sample uncertainty but also theoretical uncertainty in the galaxy model itself. Synthetic source injection provides a practical way to operationalise this: one simulation family can be injected for augmentation while another serves as a validation sample to characterise residual modelling biases, building a controlled ensemble of mismatches that can be marginalised in the redshift-calibration uncertainty budget.

\subsubsection{Methodological extensions: alternative embeddings and future inference formalisms} \label{sec:6.2.3}

As noted in Section~\ref{sec:6.1.4}, the SOM-based discretisation and the conditional density estimators adopted here have known limitations. Alternative manifold-learning methods may provide smoother neighbourhoods, improved continuity, or reduced spurious colour--redshift degeneracies in high-dimensional photometric spaces, for example diffusion maps \citep{2005math......3445N} or UMAP \citep{2018arXiv180203426M, 2025arXiv251209032A}, while latent-variable models offer a probabilistic route to continuous embeddings that avoid hard cell assignments altogether.

To make this concrete, we train a baseline UMAP embedding on the photometric \texttt{Application} datasets for both the LSST Y1 and Y10 survey scenarios, compressing the photometric feature space into two output dimensions for visualisation. The \texttt{RAIL} implementation of the SOM adopted throughout this paper operates on a compact six-feature representation consisting of one reference broad-band magnitude together with the five adjacent-band colours. UMAP imposes no analogous restriction on its input, so we take advantage of its flexibility by passing the full 21-dimensional feature vector comprising the six LSST broad-band magnitudes \(ugrizy\) together with all 15 associated pairwise colours; this is itself a practical advantage of UMAP and related neighbourhood-graph methods, which are not bound by the output-lattice trade-offs that constrain SOM design. UMAP is controlled by two principal hyperparameters: the local-neighbourhood size, which sets the scale over which the manifold topology is approximated and therefore trades attention to fine-grained local structure against preservation of the embedding's global connectivity; and the minimum inter-point distance in the embedding, which regulates how tightly adjacent points may be packed in the output space. We adopt values of \(10\) nearest neighbours and \(0.1\), respectively, together with a Euclidean input metric; all remaining hyperparameters are held at their library defaults.

Figure~\ref{fig:12} places the resulting embedding alongside the SOM discretisation adopted throughout this paper. Within each survey scenario, the upper sub-panel shows the cell-averaged true redshift on the calibration SOM grid (Section~\ref{sec:3.3}), while the lower sub-panel shows the true redshifts of individual galaxies in the UMAP two-dimensional embedding. The UMAP representation exhibits a visibly smoother colour--redshift gradient with soft transitions across the manifold, whereas the SOM shows sharper internal boundaries and more isolated high-redshift substructures at the outskirts, most clearly for LSST-Y10 where the faint, high-redshift population occupies a larger fraction of feature space. This improved smoothness reflects both UMAP's intrinsic algorithmic properties---chiefly its continuous neighbourhood-graph construction of the manifold---and the richer feature representation that it ingests. The comparison is deliberately illustrative rather than definitive: UMAP's stochastic optimisation, its continuous (rather than gridded) output space, and its reliance on post-hoc density estimation at the calibration stage would require non-trivial modifications to the importance-weighting, cluster-support, and augmented-sample-injection machinery developed in this paper. We therefore retain the SOM as the primary discretisation in the present analysis, and flag UMAP and related manifold-learning methods---together with their natural capacity to ingest richer feature representations---as a promising direction for future methodological exploration.

\begin{figure*}
    \centering
    \begin{subfigure}[t]{0.49\linewidth}
        \centering
        \includegraphics[width=\linewidth]{Figures/Y1/UMAP.pdf}
        \caption{LSST Y1}
    \end{subfigure}
    \hfill
    \begin{subfigure}[t]{0.49\linewidth}
        \centering
        \includegraphics[width=\linewidth]{Figures/Y10/UMAP.pdf}
        \caption{LSST Y10}
    \end{subfigure}
    \caption{Comparison of SOM and UMAP two-dimensional representations of the LSST photometric feature space, shown for the LSST Y1 scenario (left column) and the LSST Y10 scenario (right column). The UMAP embedding is trained with a local-neighbourhood size of \(10\), a minimum inter-point distance of \(0.1\), a Euclidean input metric, and a two-dimensional output space; see Section~\ref{sec:6.2.3} for details. Within each column, the upper sub-panel shows the cell-averaged true redshift on the calibration SOM grid (Section~\ref{sec:3.3}), while the lower sub-panel shows the true redshifts of individual galaxies in the UMAP two-dimensional embedding. UMAP yields a visibly smoother and more continuous colour--redshift manifold than the SOM, albeit without the fixed cell structure that simplifies the calibration bookkeeping adopted in this paper.}
    \label{fig:12}
\end{figure*}

On the estimation side, more robust alternatives to the \texttt{DIR} histogram include adaptive binning, \(k\)-nearest-neighbour estimators, kernel density estimators with locally varying bandwidth, mixture-density models, conditional normalising flows \citep{2019arXiv191202762P}, and hierarchical Bayesian models that borrow information across neighbouring cells. The most useful estimators will likely combine local adaptivity with explicit uncertainty propagation while remaining stable in the low-occupancy regime that drives high-redshift calibration.

More broadly, the limitations of posterior stacking discussed in Section~\ref{sec:6.1.4} motivate a shift towards fully hierarchical inference. An end-to-end Bayesian hierarchical framework can treat the population-level redshift distributions as latent objects and incorporate selection functions, transfer relations, and augmentation hypotheses directly in the generative model \citep{2019ApJ...881...80L}. The \texttt{pop-cosmos} framework \citep{2024ApJS..274...12A, 2026arXiv260203930H, 2026arXiv260203935L} exemplifies this approach, using flexible priors to infer tomographic redshift distributions jointly with calibration-layer nuisance parameters from forward-modelled galaxy populations. While such end-to-end inference is likely to be computationally demanding, the present importance-weighting framework remains valuable as a scalable approximation and diagnostic baseline.

\subsubsection{Analysis-ready products for cosmology} \label{sec:6.2.4}

For cosmological applications, the immediate goal is to produce redshift products that integrate cleanly into forecasting and inference pipelines. Even in its current form, the framework delivers: (i) an ensemble-average fiducial \(\langle \varphi_m \rangle(z)\), (ii) priors on low-dimensional nuisance parameters, and (iii) full realisations of the experiment-marginalised ensemble of tomographic redshift distributions. These outputs are already sufficient for early studies of tomographic binning, scale cuts, and robustness to alternative redshift-uncertainty models in future \(3\times2\)pt analyses.

A key next step is to quantify how much cosmological information is lost when the full uncertainty ensemble is compressed into simplified nuisance descriptions. Analytic marginalisation schemes offer an efficient route by integrating nuisance parameters out of the likelihood \citep{2020JCAP...10..056H, 2023MNRAS.522.5037R, 2023OJAp....6E..23H, 2024MNRAS.530.4412R}, but often rely on multivariate-Gaussian assumptions \citep{2010MNRAS.408..865T} that may break down for higher-order statistics. Alternatively, approaches such as \texttt{HYPERRANK} \citep{2022MNRAS.511.2170C} and basis-function expansions \citep{2021A&A...650A.148S} provide more flexible parametrisations, although the balance between parameter efficiency and representational power remains an open question. Our experiment-marginalised ensemble of redshift distributions provides a natural test-bed for these methods in the LSST context, enabling compression studies that preserve the higher-order variations most relevant to cosmological inference.

Ultimately, the objective is a unified probabilistic workflow in which spectroscopic compilation, augmentation, injection-based forward modelling, uncertainty propagation, and cosmological validation are treated consistently. The main task ahead is therefore not to replace the present methodology, but to make each component more realistic: better-motivated spectroscopic selections, broader and better-calibrated augmentation models, explicit imaging-level forward modelling, and more flexible compression of redshift-distribution uncertainty. The large suite of experiment-marginalised tomographic redshift distributions generated here is already a useful resource for forecasts, systematic comparisons of nuisance schemes, and studies of how redshift uncertainties interact with other effects, including intrinsic alignments \citep{2024PhRvD.109h3528L} and magnification bias \citep{2022MNRAS.513.1210M}.
